\documentclass[11pt]{article}
\usepackage{amsmath,amssymb,amsthm,mathtools}
\usepackage[margin=1in]{geometry}
\theoremstyle{plain}
\newtheorem{theorem}{Theorem}
\newtheorem{lemma}{Lemma}
\theoremstyle{remark}
\newtheorem{remark}{Remark}
\newcommand{\E}{\mathbb{E}}\newcommand{\PP}{\mathbb{P}}\newcommand{\R}{\mathbb{R}}
\newcommand{\eps}{\varepsilon}\newcommand{\Sch}{\mathcal{S}}
\newcommand{\ip}[2]{\langle #1,#2\rangle}

\title{Formalization brief (Mitoma campaign, 3b):\\
the Hermite--Sobolev chain generates the Schwartz topology}
\date{11 July 2026}

\begin{document}
\maketitle

\noindent\textbf{Context.} Task M3a (\texttt{mitoma3a\_report.md},
attached --- read it first) delivered
\texttt{TypeDDecouplingHermite.lean}: the Hermite functions
$h_n=c_nH_ne^{-x^2/4}$ as \texttt{hermiteBasis : HilbertBasis
$\mathbb N$ $\R$ (Lp $\R$ 2 volume)}, with \texttt{hermiteFun\_mem\_schwartz}
/ \texttt{hermiteSchwartz} and the coefficient functional on
$\Sch(\R,\R)$. This task builds the countably-Hilbertian structure on
$\Sch(\R,\R)$ that Mitoma's Lemma 3.2 (task M3c) will consume: the
Hermite--Sobolev seminorms, their equivalence with the canonical Schwartz
topology, and the explicit Hilbert--Schmidt summability between levels.
Nothing here is probabilistic; it is the last hard-analysis task of the
campaign, and most of its tools are now on the shelf.

\medskip
\noindent\textbf{Goal.} A NEW file
\texttt{TypeDDecouplingHermiteSobolev.lean} (Mathlib imports plus
\texttt{TypeDDecouplingHermite}, optionally
\texttt{...SchwartzFrechet}/\texttt{...SchwartzDual}; no existing project
file touched; whole project builds; standard axioms; no
\texttt{axiom}/\texttt{sorry}). Write $\ip{\phi}{h_n}:=\int\phi h_n$
(M3a's coefficient functional). Deliverables:

\begin{itemize}
\item[(1)] \textbf{The oscillator.} $A:\Sch(\R,\R)\to L[\R]\
\Sch(\R,\R)$, $A\phi:=-\Delta\phi+(x^2/4+1/2)\cdot\phi$, assembled from
Mathlib's \texttt{SchwartzMap.instLaplacian} (on $\R$, $\Delta=\partial^2$;
if friction arises use $-\texttt{derivCLM}\circ\texttt{derivCLM}$) and
\texttt{SchwartzMap.smulLeftCLM} (the multiplier $x\mapsto x^2/4+1/2$ has
temperate growth --- polynomial). $A$ is a continuous linear map by
construction; record also $A^r$ (iterate).
\item[(2)] \textbf{Self-adjointness on $\Sch$:}
$\int(A\phi)\psi=\int\phi(A\psi)$ for $\phi,\psi\in\Sch(\R,\R)$. The
Laplacian part is Mathlib's
\texttt{SchwartzMap.integral\_mul\_laplacian\_right\_eq\_left}; the
multiplication part is pointwise symmetry. (Backup: the zero-hypothesis
Schwartz IBP \texttt{SchwartzMap.integral\_mul\_deriv\_eq\_neg\_deriv\_mul}
applied twice.)
\item[(3)] \textbf{Eigenrelation:} $A\,h_n=(n+1)\,h_n$ (as Schwartz maps,
i.e.\ pointwise). DERIVE the eigenvalue honestly from the Hermite
identities --- our off-line computation gives, with M3a's convention
$h_n=c_nH_ne^{-x^2/4}$ and probabilists' $H_n$:
$-h_n''+(x^2/4)h_n=(n+\tfrac12)h_n$, hence eigenvalue $n+1$ for $A$; but
the paper's catch history (catch \#2, \#8) mandates checking, not
trusting. Ingredients to derive along the way (they are NOT in Mathlib;
only \texttt{Polynomial.hermite\_succ} is): the derivative identity
$H_n'=nH_{n-1}$ (induction from \texttt{hermite\_succ}), the three-term
recurrence $xH_n=H_{n+1}+nH_{n-1}$, and the resulting NORMALIZED ladder
identities, e.g.\ $x\,h_n=\sqrt{n+1}\,\alpha_{n}h_{n+1}+\sqrt n\,
\beta_nh_{n-1}$-type and $h_n'=\dots$ with explicit coefficients of size
$O(\sqrt{n+1})$ --- state them with whatever exact constants come out;
what matters downstream is the $O(\sqrt{n+1})$ size.
\item[(4)] \textbf{Coefficient decay:} for every $r\in\mathbb N$ there
are $C_r$ and a finite seminorm set with
$|\ip{\phi}{h_n}|\le(n+1)^{-r}\,\|A^r\phi\|_{L^2}\le(n+1)^{-r}C_r\cdot
(\text{finite sup of Schwartz seminorms of }\phi)$. Route: (2)+(3) give
$\ip{\phi}{h_n}=(n+1)^{-r}\ip{A^r\phi}{h_n}$; Bessel
($|\ip{\psi}{h_n}|\le\|\psi\|_{L^2}$, from M3a's basis); and
$\|\psi\|_{L^2}\le C(p_{0,0}(\psi)+p_{1,0}(\psi))$-type ($L^2$ dominated
by weighted sup norms: $\int|\psi|^2\le\sup((1+x^2)|\psi|^2)\int
(1+x^2)^{-1}$); finally continuity of the CLM $A^r$ (its seminorm bound).
\item[(5)] \textbf{Growth of Hermite seminorms:} for each $(k,m)$ there
are $C_{k,m},N_{k,m}$ with
$p_{k,m}(h_n)\le C_{k,m}(n+1)^{N_{k,m}}$ (any explicit polynomial
exponent is fine; do not chase optimality). Route: the ladder identities
(3) express $x\cdot h_n$ and $h_n'$ in $\mathrm{span}\{h_{n\pm1}\}$ with
$O(\sqrt{n+1})$ coefficients, so $x^k\partial^mh_n$ is a combination of
$h_{n-k-m},\dots,h_{n+k+m}$ with $O((n+1)^{(k+m)/2})$ coefficients;
close with the sup bound $\|h_j\|_\infty\le C(j+1)^{1/4}$ from the 1-D
Agmon inequality $\|f\|_\infty^2\le2\|f\|_{L^2}\|f'\|_{L^2}$ (prove
inline for Schwartz $f$: $f(x)^2=2\int_{-\infty}^xff'$, FTC with decay)
together with $\|h_j\|_{L^2}=1$ and $\|h_j'\|_{L^2}=O(\sqrt{j+1})$
(ladder again).
\item[(6)] \textbf{Pointwise/seminorm expansion:}
$\phi=\sum_n\ip{\phi}{h_n}h_n$ with convergence in EVERY Schwartz
seminorm (hence uniform with all derivatives and weights): by (4) with
large $r$ and (5), the tail $\sum_{n\ge N}|\ip{\phi}{h_n}|\,p_{k,m}(h_n)$
vanishes; identify the limit with $\phi$ via $L^2$ convergence (M3a's
\texttt{hermiteBasis\_hasSum\_repr}) plus continuity.
\item[(7)] \textbf{Hermite--Sobolev seminorms and the equivalence.}
$\|\phi\|_r^2:=\sum_n(n+1)^{2r}\ip{\phi}{h_n}^2$ (finite by (4)). Deliver:
(a) \texttt{sobolevSeminorm r} is a seminorm on $\Sch(\R,\R)$, HILBERTIAN
(record the inner product $\ip{\phi}{\psi}_r=\sum(n+1)^{2r}
\ip{\phi}{h_n}\ip{\psi}{h_n}$), and CONTINUOUS w.r.t.\ the canonical
topology (from (4) with $r+1$: $\|\phi\|_r^2\le C\cdot(\text{finite
seminorm sup})^2\sum(n+1)^{-2}$);
(b) the converse domination: for each $(k,m)$ there is $r$ with
$p_{k,m}(\phi)\le C\,\|\phi\|_r$ for all $\phi$ (from (5)+(6):
$p_{k,m}(\phi)\le\sum_n|\ip{\phi}{h_n}|p_{k,m}(h_n)\le\|\phi\|_r\,
(\sum_n(n+1)^{2N_{k,m}-2r})^{1/2}$, choose $r>N_{k,m}+\tfrac12$).
Together: the countable Hilbertian family $\{\|\cdot\|_r\}_{r\in\mathbb N}$
generates the Schwartz topology (two-sided lemma-level domination is the
deliverable; a \texttt{WithSeminorms} repackaging is welcome but
optional).
\item[(8)] \textbf{Hilbert--Schmidt data for M3c:} with
$e^r_j:=(j+1)^{-r}h_j$ (the $\|\cdot\|_r$-orthonormal system):
$\|e^r_j\|_q=(j+1)^{q-r}$ for $q\le r$, and the summability
$\sum_j\|e^r_j\|_q^2<\infty$ whenever $r\ge q+1$ (any sufficient
integer gap is fine; $r>q+\tfrac12$ is the sharp line). State this as a
named lemma (\texttt{hermiteSobolev\_hs\_summable}) --- it is the
nuclearity input that Mitoma's Lemma 3.2 consumes.
\end{itemize}

\begin{remark}[hygiene]
(1) DERIVE all identities and eigenvalues; the brief's stated constants
(eigenvalue $n+1$, ladder coefficient sizes, Agmon exponent $1/4$) are
predictions to check, and only the POLYNOMIAL character of the growth
matters downstream --- if a derived exponent differs, keep the honest one
and say so in the report.
(2) Fallback ladder: full delivery $\to$ (1)--(6) with (7b) for
$(k,m)=(0,0)$ only, rest documented $\to$ honest report. No fabrication.
(3) SEARCH FIRST: Bessel's inequality, \texttt{HilbertBasis} repr API,
tsum comparison/Cauchy--Schwarz for series, polynomial
\texttt{HasTemperateGrowth} instances.
(4) Report as \texttt{mitoma3b\_report.md}: the derived identities with
their exact constants, final statement names, and the (7)/(8) shapes M3c
should import.
\end{remark}

\end{document}
